\subsection{Opérations élémentaires}
\subsubsection{Mesure de l'erreur d'arrondi}
Pour $x$ et $y$ deux réels et $\star$ une opération élémentaire (addition, soustraction, multiplication, division), Pour mesurer l'erreur du calcul de $x \star y$, on note~:
$$e_\star(x,y) = \fracabs{(x + \dx)\star(y+\dy) - x \star y}{x\star y}$$
où $\abs{\delta x} \ll \abs{x}$ et $\abs{\delta y} \ll \abs{y}$.
On divise par $\abs{x \star y}$ pour avoir une mesure relative de l'erreur d'arrondi, et non absolue, c'est-à-dire dépendante du résultat de l'opération $x \star y$.
\subsubsection{Somme de deux réels}
Soient $x,y \in \R$. 
\begin{align*}
e_+(x,y) &= \frac{\abs{(x+\delta x) + (y + \delta y) - (x+y)}}{\abs{x+y}}\\
  &= \frac{\abs{\delta x + \delta y}}{\abs{x+y}}\\
  &\le \frac{\abs{\delta x}}{\abs{x+y}} + \frac{\abs{\delta y}}{\abs{x+y}}
     \qquad\text{(inégalité triangulaire)}\\
  &= \underbrace{\fracabs{\dx}{x}}_{\text{erreur initiale sur $x$}}
     \times \underbrace{\frac{\abs{x}}{\abs{x+y}}}_{\text{coefficient amplificateur}}
   + \underbrace{\fracabs{\dy}{y}}_{\text{erreur initiale sur $y$}}
     \times \underbrace{\frac{\abs{y}}{\abs{x+y}}}_{\text{coefficient amplificateur}}
\end{align*}

\begin{remarque}{}{rôle des coefficients amplificateurs}
    L’erreur relative sur la somme est une combinaison linéaire des erreurs relatives
    initiales sur $x$ et $y$, pondérées par les coefficients
    \[
    \frac{|x|}{|x+y|} \quad \text{et} \quad \frac{|y|}{|x+y|}.
    \]
    Lorsque $x$ et $y$ sont de même signe et de tailles comparables, ces coefficients
    sont de l’ordre de $1$, et l’erreur est peu amplifiée.

    En revanche, lorsque $x$ et $y$ sont de signes opposés et que $|x+y| \ll |x|, |y|$,
    les coefficients amplificateurs deviennent très grands : de petites erreurs
    initiales sur $x$ et $y$ peuvent alors produire une grande erreur relative sur
    la somme. Ce phénomène correspond à une perte de chiffres significatifs
    (annulation catastrophique, voir \href
{https://fr.wikipedia.org/wiki/Annulation_catastrophique}
{Wikipédia}).
\end{remarque}

\begin{remarque}{}{erreurs initiales}
    Les formules des erreurs initiales $\fracabs{\dx}{x}$ et $\fracabs{\dy}{y}$, proviennent de l'expression~:

    $$\fracabs{(x + \delta x) - x}{x} = \fracabs{\dx}{x} $$
    qui mesure l'erreur entre la valeur exacte $x$ et la valeur approchée $x + \delta x$. Là encore on divise par $\abs{x}$ pour avoir une mesure relative de l'erreur, et non absolue, c'est-à-dire dépendante de la valeur exacte $\abs{x}$.
    
\end{remarque}

\subsubsection{Produit de deux réels}
Soient $x,y \in \R$. 
\begin{align*}
e_\times(x,y) &= \frac{\abs{(x+\delta x) \times (y + \delta y) - x \times y}}{\abs{x \times y}}\\
  &= \frac{\abs{x \delta y + y \delta x + \delta x \delta y}}{\abs{x \times y}}\\
  &\le \frac{\abs{x \delta y}}{\abs{x \times y}} + \frac{\abs{y \delta x}}{\abs{x \times y}} + \frac{\abs{\delta x \delta y}}{\abs{x \times y}}\\
  &= \underbrace{\fracabs{\dy}{y}}_{\text{erreur initiale sur $y$}}
     + \underbrace{\fracabs{\dx}{x}}_{\text{erreur initiale sur $x$}}
     + \underbrace{\fracabs{\dx}{x} \times \fracabs{\dy}{y}}_{\text{terme négligeable}}
\end{align*}
\begin{remarque}{}{sur l'expression}
    L'erreur relative sur le produit n'est que la somme des erreurs relatives initiales sur $x$ et $y$, en plus d'un terme négligeable du second ordre, qu'on peut négliger dans un cadre simple.\\
    Dans la mesure où les ordres de grandeurs s'additionnent également dans un produit, on peut dire que le produit est une opération stable du point de vue des erreurs d'arrondi.
\end{remarque}

\subsection{Principes généraux pour réduire les erreurs d'arrondi}

\subsubsection{Éviter de soustraire deux quantités absoluements proches}

\begin{exemple}{}{soustraction au cosinus}
    POur calculer $y = 1-\cos(x)$, pour $x$ proche de 0, on peut utiliser la formule~:
    $$\cos(2x) = 1 - 2\sin^2(x)$$
    pour obtenir~:
    $$y = 2 \sin^2(x/2)$$
\end{exemple}

\begin{exemple}{}{}
    On s'intéresse à la recherche des racines de $P = aX^2 + bX + c$, de discriminant $\Delta$ vérifiant~:
    $$\Delta \gg 0$$
    On a alors~:
    $$b^2 \gg 4ac$$
    Ce discriminant est donc proche de $b^2$. Ainsi, quand on calcule~:
    $$x_1 = \frac{-b-\Delta}{2a} \qquad \text{et} \qquad x_2 = \frac{-b + \Delta}{2a}$$ 
    le calcul de $x2$ est imprécis à cause du dénominateur. Mais via les formules de Viète, $x_2 = \frac{c}{ax_1}$.
\end{exemple}

\subsubsection{Sommer d'abord les termes les plus petits}
Sommer un grand nombre avec un petit nombre entraîne la troncature des plus petites décimales du petit nombre.
\begin{exemple}{}{calcul d'une somme}
    Pour calculer~:
    $$S = \sum_{i=1}^n \frac{1}{i}$$
    bien qu'il soit naturel de sommer les termes dans l'ordre croissant des indices, l'ordre rétrograde est ici à privilégier. On calcule sur quatre décimales~:\\
    \begin{center}
        \begin{tabular}{c|c|c|c}
            $n$ & croissant & rétrograde & exact\\
            \hline
            $100$ & $5,187$ & $5,187$ & $5,1873775$\\
            $1\,000$ & $7,449$ & $7,485$ & $7,485471$\\
            $10\,000$ & $8,449$ & $9,448$ & $9,787606$\\
            
        \end{tabular}
    \end{center}
\end{exemple}
